\documentclass[a4paper]{article}

\usepackage[margin=1in]{geometry}
\usepackage{amsmath,amsfonts,mathtools,amsthm,amssymb}
\usepackage[l2tabu,orthodox]{nag}
\usepackage{microtype} % fixes spacing or whatever
\usepackage{enumitem}
\usepackage{tikz-cd}
\usepackage{xcolor}
\usepackage{physics}
\usepackage{mathrsfs}
\usepackage{cancel}
% \usepackage{libertine,libertinust1math}
\usepackage{eucal}
\usepackage[toc,page]{appendix}
\usepackage{pgfplots}
\pgfplotsset{compat=1.18}

\newtheorem{proposition}{Proposition}
\newtheorem{theorem}{Theorem}
\newtheorem{lemma}{Lemma}
\newtheorem{corollary}{Corollary}

\theoremstyle{definition}
\newtheorem{remark}{Remark}
\newtheorem{definition}{Definition}
\newtheorem{example}{Example}

% \usepackage[sc,osf]{mathpazo}   % With old-style figures and real smallcaps.
% \linespread{1.025}              % Palatino leads a little more leading
% % Euler for math and numbers
% \usepackage[euler-digits,small]{eulervm}
% \AtBeginDocument{\renewcommand{\hbar}{\hslash}}

\usepackage{etoolbox}
\usepackage{dynkin-diagrams}
\def\row#1/#2!{\dynkin{#1}{#2}\\}
\newcommand{\tble}[1]{
   \renewcommand*\do[1]{\row##1!}
   \[
      \begin{array}{ll}\docsvlist{#1}\end{array}
   \]
}

\allowdisplaybreaks

\renewcommand{\epsilon}{\varepsilon}
% \renewcommand{\phi}{\varphi}

\DeclareMathAlphabet{\mathpzc}{OT1}{pzc}{m}{it}

\newcommand{\goth}[1]{\mathfrak{#1}}
\newcommand{\red}[1]{#1_{\text{red}}}
\newcommand{\ad}[1]{#1_{\text{Ad}}}
\newcommand{\reg}[1]{#1_{\text{reg}}}
\newcommand{\sh}[1]{#1^{\text{sh}}}
\newcommand{\cpdv}[2]{\cfrac{\partial #1}{\partial #2}}

\newcommand{\fA}{\mathcal{A}}
\newcommand{\fB}{\mathcal{B}}
\newcommand{\fC}{\mathcal{C}}
\newcommand{\fD}{\mathcal{D}}
\newcommand{\fE}{\mathcal{E}}
\newcommand{\fF}{\mathcal{F}}
\newcommand{\fG}{\mathcal{G}}
\newcommand{\fH}{\mathcal{H}}
\newcommand{\fI}{\mathcal{I}}
\newcommand{\fJ}{\mathcal{J}}
\newcommand{\fK}{\mathcal{K}}
\newcommand{\fL}{\mathcal{L}}
\newcommand{\fM}{\mathcal{M}}
\newcommand{\fN}{\mathcal{N}}
\newcommand{\fO}{\mathcal{O}}
\newcommand{\fP}{\mathcal{P}}
\newcommand{\fQ}{\mathcal{Q}}
\newcommand{\fR}{\mathcal{R}}
\newcommand{\fS}{\mathcal{S}}
\newcommand{\fT}{\mathcal{T}}
\newcommand{\fX}{\mathcal{X}}
\newcommand{\fY}{\mathcal{Y}}
\newcommand{\fZ}{\mathcal{Z}}

\newcommand{\be}{\mathbf{e}}
\newcommand{\bx}{\mathbf{x}}
\newcommand{\bone}{\mathbf{1}}
\newcommand{\cx}{\bar{x}}
\newcommand{\cy}{\bar{y}}
\newcommand{\vx}{\vec{x}}
\newcommand{\vsigma}{\vec{\sigma}}
\newcommand{\tzeta}{\tilde{\zeta}}

\newcommand{\gm}{\goth{m}}
\newcommand{\gp}{\goth{p}}
\newcommand{\gF}{\goth{F}}
\newcommand{\gU}{\goth{U}}
\newcommand{\gV}{\goth{V}}

\newcommand{\A}{\mathbf{A}}
\newcommand{\C}{\mathbf{C}}
\newcommand{\F}{\mathbf{F}}
\newcommand{\G}{\mathbf{G}}
\newcommand{\bH}{\mathbf{H}}
\newcommand{\N}{\mathbf{N}}
\newcommand{\PP}{\mathbf{P}}
\newcommand{\Q}{\mathbf{Q}}
\newcommand{\R}{\mathbf{R}}
\newcommand{\Z}{\mathbf{Z}}

\newcommand{\dlog}{\dd\log}

\newcommand\srestr[2]{{\left.\kern-\nulldelimiterspace #1\vphantom{\small|} \right|_{#2}}}
\newcommand\restr[2]{{\left.\kern-\nulldelimiterspace #1 \vphantom{\big|} \right|_{#2}}}
\newcommand\gsec[2]{\Gamma({#1}, {#2})}
\newcommand\gsecx[1]{\Gamma(X, {#1})}

\DeclareMathOperator{\chr}{char}
\DeclareMathOperator{\rProj}{\mathbf{Proj}}
\DeclareMathOperator{\rSpec}{\mathbf{Spec}}
\DeclareMathOperator{\rHom}{\mathbf{Hom}}
\DeclareMathOperator{\rExt}{\mathbf{Ext}}
\DeclareMathOperator{\id}{id}
\DeclareMathOperator{\Frac}{Frac}
\DeclareMathOperator{\rk}{rank}
\DeclareMathOperator{\deter}{det}
\DeclareMathOperator{\pic}{Pic}
\DeclareMathOperator{\cacl}{CaCl}
\DeclareMathOperator{\trd}{tr.d.}
\DeclareMathOperator{\cl}{Cl}
\DeclareMathOperator{\depth}{depth}
\DeclareMathOperator{\codim}{codim}
\DeclareMathOperator{\Div}{Div}
\DeclareMathOperator{\coker}{coker}
\DeclareMathOperator{\len}{length}
\DeclareMathOperator{\height}{height}
\DeclareMathOperator{\supp}{Supp}
\DeclareMathOperator{\proj}{Proj}
\DeclareMathOperator{\im}{im}
\DeclareMathOperator{\Hom}{Hom}
\DeclareMathOperator{\Der}{Der}
\DeclareMathOperator{\spec}{Spec}
\DeclareMathOperator{\Aut}{Aut}
\DeclareMathOperator{\ch}{char}
\DeclareMathOperator{\tor}{Tor}
\DeclareMathOperator{\Ann}{Ann}
\DeclareMathOperator{\Syl}{Syl}
\DeclareMathOperator{\Sym}{Sym}
\DeclareMathOperator{\GL}{GL}
\DeclareMathOperator{\SL}{SL}
\DeclareMathOperator{\Stab}{Stab}
\DeclareMathOperator{\Perm}{Perm}
\DeclareMathOperator{\Orb}{Orb}
\DeclareMathOperator{\Gal}{Gal}
\DeclareMathOperator{\Supp}{Supp}
\DeclareMathOperator{\Ext}{Ext}
\DeclareMathOperator{\Tor}{Tor}
\DeclareMathOperator{\PGL}{PGL}
\DeclareMathOperator{\hd}{hd}
\DeclareMathOperator{\pd}{pd}
\DeclareMathOperator{\SO}{SO}
\DeclareMathOperator{\SU}{SU}
\DeclareMathOperator{\Sp}{Sp}
\DeclareMathOperator{\Oth}{O}
\DeclareMathOperator{\Uni}{U}
\DeclareMathOperator{\evf}{ev}
\DeclareMathOperator{\cH}{H}
\DeclareMathOperator{\dR}{R}
\DeclareMathOperator{\End}{End}
\DeclareMathOperator{\Hasse}{Hasse}
\DeclareMathOperator{\Num}{Num}

\author{James Lee}

% Solarized
% \pagecolor[RGB]{0,20,26}
% \color[RGB]{191,191,191}

% Sephia
% \pagecolor[RGB]{249,239,220}

% Grey
% \pagecolor[RGB]{200, 200, 200}

% Black
% \pagecolor[RGB]{0,0,0}
% \color[RGB]{255,255,255}


\title{Chapter 5, Section 1}

\begin{document}

\maketitle

\begin{enumerate} [label=\textbf{\arabic*.}, leftmargin=0em]

    % 1.1, 1.2, 1.3, 1.4, 1.5, 1.8, 1.10, 2.1, 2.2, 2.3, 2.6, 2.8, 2.16, 3.1, 3.5, 3.7, 4.3, 4.4, 4.5, 4.12, 4.16, 5.1, 5.5, 5.6, 6.1

    \item Let $C, D$ be any two divisors on a surface $X$, and let the corresponding invertible sheaves be $\fL$ and $\fM$.
          Show that
          \begin{equation*}
              C.D = \chi(\fO_X) - \chi(\fL^{-1}) - \chi(\fM^{-1}) + \chi(\fL^{-1} \otimes \fM^{-1}).
          \end{equation*}

          \begin{proof}
              Applying the Riemann-Roch theorem to the divisor $-C - D$ gives
              \begin{equation*}
                  \chi(\fO_X(-C - D)) = \frac{1}{2}(-C - D).(-C - D - K) + \chi(\fO_X)
              \end{equation*}
              The first term on the right hand side simplies to
              \begin{equation*}
                  (-C - D).(-C - D - K) = C^2 + 2C.D. + D^2 + C.K + D.K.
              \end{equation*}
              Substitutting back to above and isolating $C.D$ gives
              \begin{align*}
                  C.D & = -\chi(\fO_X) + \chi(\fO_X(-C - D)) - \frac{1}{2}C.(C + K) - \frac{1}{2}C.(C + K)                                                   \\
                      & = \chi(\fO_X) + \chi(\fO_X(-C - D)) - \bigg(\chi(\fO_X) - \frac{1}{2}C.(-C-K) \bigg) - \bigg(\chi(\fO_X) -\frac{1}{2}D.(-D -K)\bigg) \\
                      & = \chi(\fO_X) + \chi(\fO_X(-C-D)) - \chi(\fO_X(-C)) - \chi(\fO_X(-D)).
              \end{align*}
          \end{proof}

    \item Let $H$ be a very ample divisor on the surface $X$, corresponding to a projective embedding $X \subseteq \P^N$.
          If we write the Hilbert polynomial of $X$ (III, Ex. 5.2) as
          \begin{equation*}
              P(z) = \frac{1}{2}az^2 + bz + c,
          \end{equation*}
          show that $a = H^2$, $b = \frac{1}{2}H^2 + 1 - \pi$, where $\pi$ is the genus of a nonsingular curve representing $H$, and $c = 1 + p_a$.
          Thus the \textit{degree} of $X$ in $\P^N$, as defined in (I,  \SS 7), is just $H^2$.
          Show also that if $C$ is any curve in $X$, then the degree of $C$ in $\P^N$ is just $C.H$.

          \begin{proof}
              By the Riemann-Roch theorem, we have the equality
              \begin{equation*}
                  P(z) = \frac{1}{2}H^2z^2 - \frac{1}{2}(H.K)z + 1 + p_a.
              \end{equation*}
              Matching coefficients readily gives $a, b$, and $c$.
              We can rewrite the coefficient of $z$ as
              \begin{equation*}
                  -\frac{1}{2}(H.K) = \frac{1}{2}H^2 - \frac{1}{2} H.(H + K),
              \end{equation*}
              and the adjunction formula (1.5) gives
              \begin{equation*}
                  2\pi - 2 = H.(H + K).
              \end{equation*}
              We conclude that
              \begin{equation*}
                  b = \frac{1}{2}H^2 + 1 - \pi.
              \end{equation*}
              If $C$ is any curve in $X$, then the degree of $C$ in $\P^N$ is just the degree of $\fO_X(H) \otimes \fO_C$ as a line bundle on $C$, which is equal to $C.H$ by (1.3).
          \end{proof}

    \item Recall that the \textit{arithmetic genus} of a projective scheme $D$ of dimension 1 is defined as $p_a(D) = 1 - \chi(\fO_D)$ (III, Ex. 5.3).
          \begin{itemize}
              \item[(a)] If $D$ is an effective divisor on a surface $X$, use (1.6) to show that $2p_a(D) - 2 = D.(D + K)$.
              \item[(b)] $p_a(D)$ depends only on the linear equivalence class of $D$ on $X$.
              \item[(c)] More generally, for \textit{any} divisor $D$ on $X$, we define the \textit{virtual arithmetic genus} (which is equal to the ordinary arithmetic genus if $D$ is effective) by the same formula: $2p_a - 2 = D.(D + K)$.
                    Show that for any two divisors $C, D$, we have
                    \begin{equation*}
                        p_a(-D) = D^2 - p_a(D) + 2,
                    \end{equation*}
                    and $$p_a(C + D) = p_a(C) + p_a(D) + C.D - 1.$$
          \end{itemize}

          \begin{proof} $ $ \vspace{0pt}
              \begin{itemize} [leftmargin=0cm]
                  \item[(a)] We have the exact sequence
                        \[ \begin{tikzcd}
                                0 \arrow[r] & \fO_X(-D) \arrow[r] & \fO_X \arrow[r] & \fO_D \arrow[r] & 0,
                            \end{tikzcd} \]
                        and taking Euler characteristics gives
                        \begin{equation*}
                            \chi(\fO_D) = \chi(\fO_X) - \chi(\fO_X(-D)).
                        \end{equation*}
                        On the other hand, the Riemann-Roch theorem gives
                        \begin{equation*}
                            \chi(\fO_X(-D)) = \frac{1}{2}D.(D + K) + \chi(\fO_X).
                        \end{equation*}
                        Substituting gives
                        \begin{equation*}
                            1 - \chi(\fO_D) = 1 + \frac{1}{2}D.(D + K).
                        \end{equation*}

                  \item[(b)] Follows from (a) and (1.1(4)), which states $D.(D + K)$ depends only on the linear equivalence class of $D$ on $X$.
                  \item[(c)]
                        \begin{align*}
                            2p_a(-D) - 2    & = -D.(-D + K)                        \\
                                            & = D.(D - K)                          \\
                                            & = 2D^2 - D.(D + K)                   \\
                                            & = 2D^2 - 2p_a(D) + 2                 \\
                            2p_a(C + D) - 2 & = (C + D).(C + D + K)                \\
                                            & = C.(C + K) + D.(D + K) + 2(C.D)     \\
                                            & = 2p_a(C) - 2 + 2p_a(D) - 2 + 2(C.D)
                        \end{align*}
              \end{itemize}
          \end{proof}

    \item \begin{itemize}
              \item[(a)] If a surface $X$ of degree $d$ in $\P^3$ contains a straight line $C = \P^1$, show that $C^2 = 2 - d$.
              \item[(b)] Assume $\car{k} = 0$, and show for every $d \geq 1$, there exists a nonsingular surface $X$ of degree $d$ in $\P^3$ containing the line $x = y = 0$.
          \end{itemize}

          \begin{proof} $ $ \vspace{0pt}
              \begin{itemize} [leftmargin=0cm]
                  \item[(a)] By (Ex. 1.5), $C^2 = -C.K - 2$.
                        By (1.3), $C.K = d - 4$, because $C$ is a straight line and $\fO_X(K) = \fO_X(d - 4)$.
                        Hence, $C^2 = 2 - d$.

                  \item[(b)] We can prove the following more general statement:
                        \begin{proposition}
                            For every $n \geq 3$ and $d > 0$, and for every linear subspace $H \subset \P^n$ of codimension at least 2, there exists a nonsingular hypersurface of degree $d$ in $\P^n$ containing $H$.
                        \end{proposition}
                        We can assume $H$ has codimension two and is defined by the vanishing of $x = y = 0$.
                        Bertini's theorem (III, 10.9.2) says a general hypersurface $X$ of degree $d$ in $\P^3$ containing $H$ is nonsingular away from $H$.
                        Also, any such $X$ is is defined by the vanishing of $xF + yG = 0$, where $F, G$ define surfaces of degree $d - 1$.
                        If $P$ is a point in $H$, then $X$ is singular at $P$ if and only if both $F, G$ vanish at $P$.
                        This is a closed condition, so we can find plenty of choices of $F, G$ which give a nonsingular hypersurface containing $H$.

                        \begin{example}
                            For $n = 3$ and $d \geq 2$, the surface defined by $$x(z^{d - 1} + x^{d - 1}) + y(w^{d - 1} + y^{d - 1}) = 0$$ is nonsingular and contains $L$.
                        \end{example}
              \end{itemize}
          \end{proof}

    \item \begin{itemize}
              \item[(a)] If $X$ is a surface of degree $d$ in $\P^3$, then $K^2 = d(d - 4)^2$.
              \item[(b)] If $X$ is a product of two nonsingular curves $C, C'$, of genus $g, g'$ respectively, then $K^2 = 8(g - 1)(g' - 1)$.
          \end{itemize}

          \begin{proof} $ $ \vspace{0pt}
              \begin{itemize} [leftmargin=0cm]
                  \item[(a)] The adjunction formula gives $\fO_X(K) = \fO_X(d - 4)$.
                        If $H$ is a general hyperplane section of $X$ in $\P^3$, then $K = (d - 4)H$, and $H^2 = d$ by (Ex. 1.2).
                        Hence, $$K^2 = (d - 4)^2H^2 = d(d - 4)^2.$$

                  \item[(b)] By (II, Ex. 8.3), there is an isomorphism $\omega_{X} \cong p_1^*\omega_C \otimes p_2^*\omega_{C'}$.
                        In particular, if $\gk'$ and $\gk'$ are canonical divisors of $C$ and $C'$, respectively, then $p_1^*\gk + p_2^*\gk'$ is a canonical divisor of $X$.
                        Denote $D = p_1^*\gk$ and $D' = p_2^*\gk'$.
                        The support of $D$ consists of $2g - 2$ copies of $C'$ which are the fibres of the points in $\gk$ with respect to $p_1 : X \to C$.
                        It follows that
                        \begin{equation*}
                            D.D' = (\deg{K_C}) \cdot (\deg{K_{C'}}) = 4(g - 1)(g' - 1).
                        \end{equation*}
                        Also, we can always find two canonical divisors with disjoint support, and since the intersection pairing only depends on the linear equivalence class, we have $D^2 = D'^2 = 0$.
                        We conclude that $K^2 = 2D.D' = 8(g - 1)(g' - 1)$.
              \end{itemize}
          \end{proof}

    \item \begin{itemize}
              \item[(a)] If $C$ is a curve of genus $g$, show that the diagonal $\Delta \subseteq C \times C$ has self-intersection $\Delta^2 = 2 - 2g$. (Use the definition of $\omega_C$ in (II, \SS 8).)
              \item[(b)] Let $l = C \times \text{pt}$ and $m = \text{pt} \times C$.
                    If $g \geq 1$, show that $l$, $m$, and $\Delta$ are linearly independent in $\Num(C \times C)$.
                    Thus $\Num(C \times C)$ has rank $\geq 3$, and in particular, $\pic(C \times C) \neq p_1^*\pic{C} \oplus p_2^*\pic{C}$.
          \end{itemize}

          \begin{proof} $ $ \vspace{0pt}
              \begin{itemize} [leftmargin=0cm]
                  \item[(a)] By definition, $\omega_{C} \cong \fO_X(-\Delta) \otimes \fO_\Delta$, and by (1.3), $\deg_\Delta(\fO_X(-\Delta) \otimes \fO_\Delta) = -\Delta^2$.

                  \item[(b)] Suppose $D = a l + b m + c \Delta = 0$ in $\Num(C \times C)$ for some integers $a, b, c \in \Z$.
                        It is not hard to see
                        \begin{equation*}
                            l^2 = m^2 = 0, \quad l.m = 1, \quad \Delta.l = \Delta.m = 1,
                        \end{equation*}
                        which implies
                        \begin{equation*}
                            D.m = b + c = 0, \quad D.l = a + c = 0, \quad D.\Delta = a + b + c(2 - 2g) = 0.
                        \end{equation*}
                        Combining the first two gives $a + b + 2c = 0$.
                        The condition $g \geq 1$ forces $a, b, c = 0$, which is what we wanted to show.
              \end{itemize}
          \end{proof}

    \item \textit{Algebraic Equivalence of Divisors.} Let $X$ be a surface.
          Recall that we have defined an algebraic family of effective divisors on $X$, parametrized by a nonsingular curve $T$, to be an effective Cartier divisor $D$ on $X \times T$, flat over $T$ (III, 9.8.5).
          In this case, for any two closed points $0, 1 \in T$, we say that the corresponding divisors $D_0, D_1$ on $X$ are prealgebraically equivalent.
          Two arbitrary divisors are prealgebraically equivalent if they are differences of prealgebraically equivalent divisors.
          Two divisors $D, D'$ are \textit{algebraically equivalent} if there is a finite sequence $D = D_0, D_1, \dots, D_n = D'$ with $D_i$ and $D_{i + 1}$ prealgebraically equivalent for each $i$.
          \begin{itemize}
              \item[(a)] Show that the divisors algebraically equivalent to 0 form a subgroup of $\Div{X}$.
              \item[(b)] Show that linearly equivalent divisors are algebraically equivalent.
              \item[(c)] Show that algebraically equivalent divisors are numerically equivalent.
          \end{itemize}

          \begin{proof} $ $ \vspace{0pt}
              \begin{itemize} [leftmargin=0cm]
                  \item[(a)]

                  \item[(b)]

                  \item[(c)]
              \end{itemize}
          \end{proof}

    \item \textit{Cohomology Class of a Divisor.} For any divisor $D$ on the surface $X$, we define its cohomology class $c(D) \in \cH^1(X, \Omega_X)$ by using the isomorphism $\pic{X} \cong \cH^1(X, \fO_X^*)$ of (III, Ex. 4.5) and the sheaf homomorphism $\dlog : \fO_X^* \to \Omega_X$ (III, Ex. 7.4c).
          Thus we obtain a group homomorphism $c : \pic{X} \to \cH^1(X, \Omega_X)$.
          On the other hand, $\cH^1(X, \Omega_X)$ is dual to itself by Serre duality (III, 7.13), so we have a nondegenerate bilinear map
          \begin{equation*}
              \langle \quad, \quad \rangle : \cH^1(X, \Omega_X) \times \cH^1(X, \Omega_X) \to k.
          \end{equation*}
          \begin{itemize}
              \item[(a)] Prove that this is compatible with the intersection pairing, in the following sense: for any two divisors $D, E$ on $X$, we have
                    \begin{equation*}
                        \langle c(D), c(E) \rangle = (D.E) \cdot 1
                    \end{equation*}
                    in $k$.
              \item[(b)] If $\car{k} = 0$, use the fact that $\cH^1(X, \Omega_X)$ is a finite-dimensional vector space to show that $\Num{X}$ is a finitely generated free abelian group.
          \end{itemize}

          \begin{proof} $ $ \vspace{0pt}
              \begin{itemize} [leftmargin=0cm]
                  \item[(a)] We need to first show numerically equivalent divisors give isomorphic cohomology classes.
                        Suppose $D \equiv 0$ and $D \neq 0$ in $\pic{X}$.
                        Let $(U, f)$ be a locally defining equation of $D$.
                        We want to show $\dd f / f = 0$.
                        Let $P$ be any point in $X$,

                  \item[(b)]
              \end{itemize}
          \end{proof}

    \item \begin{itemize}
              \item[(a)] If $H$ is an ample divisor on the surface $X$, and if $D$ is any divisor, show that
                    \begin{equation*}
                        (D^2)(H^2) \leq (D.H)^2.
                    \end{equation*}
              \item[(b)] Now let $X$ be a product of two curves $X = C \times C'$.
                    Let $l = C \times \text{pt}$, and $m = \text{pt} \times C'$.
                    For any divisor $D$ on $X$, let $a = D.l$, $b = D.m$.
                    Then we say $D$ has \textit{type} $(a, b)$.
                    If $D$ has type $(a, b)$, with $a, b \in \Z$, show that
                    \begin{equation*}
                        D^2 \leq 2ab,
                    \end{equation*}
                    and equality holds if and only if $D = bl + am$.
          \end{itemize}

          \begin{proof} $ $ \vspace{0pt}
              \begin{itemize} [leftmargin=0cm]
                  \item[(a)] By replacing $H$ with $mH$ for some $m \ag 0$ if necessary, we may assume that $H$ is an irreducible nonsingular curve corresponding to a very ample line bundle on $X$.
                        Consider the case when $D$ is an irreducible curve.

                  \item[(b)]
              \end{itemize}
          \end{proof}

          \begin{example}
              Consider the quadric surface $Q = \P^1 \times \P^1$ and the divisor $H$ of type $(1, 1)$.
              If $D$ be a divisor of type $(a, b)$, then by (1.4.3), the left-hand side is equal to $2ab$, while the right-hand side is equal to $(a + b)^2$ (IV, 6.4.1) (Ex. 1).
          \end{example}

    \item \textit{Weil's Proof [2] of the Analogue of the Riemann Hypothesis for Curves.} Let $C$ be a curve of genus $g$ defined over the finite field $\F_q$, and let $N$ be the number of points of $C$ rational over $\F_q$.
          Then $N = 1 - a + q$, with $|a| \leq 2g \sqrt{q}$.
          To prove this, we consider $C$ as a curve over the algebraic closure $k$ of $\F_q$.
          Let $f : C \to C$ be the $k$-linear Frobenius morphism obtained by taking $q$th powers, which makes sense since $C$ is defined over $\F_q$, so $X_q \cong X$ (IV, 2.4.1).
          Let $\Gamma \subseteq C \times C$ be the graph of $f$, and let $\Delta \subseteq C \times C$ be the diagonal.
          Show that $\Gamma^2 = q(2 - 2g)$, and $\Gamma.\Delta = N$.
          Then apply (Ex. 1.9) to $D = r\Gamma + s\Delta$ for all $r$ and $s$ to obtain the result.

          \begin{proof}
              
            For any non-zero $r, s \in \Z$, consider the divisor $D = r\Gamma - s\Delta$.
            It has type $(r - s, r - s)$, so by (Ex. 1.9), we have the inequality
            \begin{equation*}
                D^2 \leq 2(r - s)^2.
            \end{equation*}
          \end{proof}

    \item In this problem, we assume that $X$ is a surface for which $\Num{X}$ is finitely generated.
          \begin{itemize}
              \item[(a)] If $H$ is an ample divisor on $X$, and $d \in \Z$, show that the set of effective divisors $D$ with $D.H = d$, modulo numerical equivalence, is a finite set.
              \item[(b)] Now let $C$ be a curve of genus $g \geq 2$, and use (a) to show that the group of automorphisms of $C$ is finite, as follows.
                    Given an automorphism $\sigma$ of $C$, let $\Gamma \subseteq X = C \times C$ be its graph.
                    First show that if $\Gamma \equiv \Delta$, then $\Gamma = \Delta$, using the fact that $\Delta^2 < 0$, since $g \geq 2$ (Ex. 1.6).
                    Then use (a).
          \end{itemize}

    \item If $D$ is an ample divisor on the surface $X$, and $D' \equiv D$, then $D'$ is also ample.
          Give an example to show, however, that if $D$ is very ample, $D'$ need not be very ample.

          \begin{proof}

          \end{proof}

\end{enumerate}

\end{document}
